\documentclass{article}
\usepackage{amsmath, amssymb, amsthm}

\setlength{\parindent}{0pt}

\newtheorem*{claim}{Claim}

\begin{document}

\begin{claim}
    If the equation $A\mathbf{x} = \mathbf{0}$ has a unique solution, then $A$ is left invertible.
\end{claim}

\begin{proof}
    Let $A$ be an $m \times n$ matrix. Assume that the equation $A\mathbf{x} = \mathbf{0}$ has a unique 
    solution, so there are no free variables, i.e. $A$ has a pivot in every column of its reduced echelon form. 
    Let $A_e$ be the reduced echelon form of $A$, so 
    \[
    A_e = EA
    \]
    where $E$ is a product of elementary matrices, corresponding to the row operations that reduce $A$. Since 
    every column of $A_e$ has a pivot, there are $n$ pivots. By definition of a pivot, there can be at most 
    one pivot in each row. Therefore $n \leq m$, and 
    \[
    A_e = EA = 
    \begin{bmatrix}
        I_n \\ \mathbf{0}_{(m-n) \times n}
    \end{bmatrix}
    \]
    where $I_n$ is the $n \times n$ identity matrix and $\mathbf{0}_{(m-n) \times n}$ is the $(m-n) \times n$ 
    zero matrix. Define the $n \times m$ matrix 
    \[
    P = 
    \begin{bmatrix}
        I_n & \mathbf{0}_{n \times (m-n)}
    \end{bmatrix}
    \]
    where $\mathbf{0}_{n \times (m-n)}$ is the $n \times (m-n)$ zero matrix. Then
    \[
    P(EA) = P\begin{bmatrix} I_n \\ \mathbf{0}_{(m-n) \times n} \end{bmatrix} = I_n
    \]
    By associativity, 
    \[
    (PE)A = P(EA) = I_n
    \]
    Thus $PE$ is a left inverse of $A$, so $A$ is left invertible.

    \medskip

    Hence if the equation $A\mathbf{x} = \mathbf{0}$ has a unique solution, then $A$ is left invertible.
\end{proof}

\end{document}
